\documentclass[11pt, oneside]{article} 
\usepackage{geometry}
\geometry{letterpaper} 
\usepackage{graphicx}
	
\usepackage{amssymb}
\usepackage{amsmath}
\usepackage{parskip}
\usepackage{color}
\usepackage{hyperref}

\graphicspath{{/Users/telliott/Github-Math/figures/}}

\title{Perpendicular bisector and SSS}
\date{}

\begin{document}
\maketitle
\Large

%[my-super-duper-separator]

\subsection*{perpendicular bisector}

Suppose are given that $BD = DC$ and that $\angle ADC$ is a right angle.  For any point on $AD$, we can draw triangles with vertices at $B$ and $C$ that are congruent by SAS.

\begin{center} \includegraphics [scale=0.4] {iso13.png} \end{center}

We conclude that every point on $AD$ (and extensions of $AD$) is equidistant from $B$ and $C$.

The converse theorem says that \emph{every} point which is equidistant from $B$ and $C$ lies on $AD$.  

Here is a quick proof relying on the \hyperref[sec:triangle_inequality]{\textbf{triangle inequality}}, which says that in any triangle, the sum of any two sides must be greater than the length of the third side.  We will also prove this one later.  To the current proof:

\emph{Proof}.

Suppose that $P$ is equidistant from $B$ and $C$ but does not lie on the perpendicular bisector.  Then, erect the perpendicular bisector and find the point where $PC$ crosses the bisector at $A$.
\begin{center} \includegraphics [scale=0.4] {iso13c.png} \end{center}
By the forward theorem, $AB = AC$.

We are supposing that $PB = PC$.  By the triangle inequality
\[ PB < AB + AP \]
Since $AB = AC$:
\[ PB < AC + AP = PC \]
But this is absurd.  $PB$ cannot both be equal to and less than $PC$.  Therefore, our supposition is incorrect, and there does not exist any such point $P$.

$\square$

See \hyperref[sec:perp_bi_converse]{\textbf{here}} for a somewhat fuller explanation.

\subsection*{SSS $\rightarrow$ SAS}

\label{sec:SSS_implies_SAS}

We will prove that the SSS criterion implies SAS.

\emph{Proof}.

Let all the sides of $\triangle ABC$ be equal to $\triangle DEF$.  Choose the longest side of $\triangle ABC$ and align it with $\triangle DEF$ as shown.  Since $AB$ is coincident with $DE$ we suppress the latter labels.  $\triangle DEF$ is drawn as the mirror image of $\triangle ABC$.  That is allowed for congruency, and if $\triangle DEF$ were directly superimposable, we could just flip it.

\begin{center} \includegraphics [scale=0.4] {SSS.png} \end{center}

$D$ is placed so that two sets of sides are equal, and equal sides are adjacent in the quadrilateral.  $AC = AD$ and $BC = BD$.  Of course, the third side $AB$ is shared.  So we have SSS.

By the forward version of the isosceles triangle theorem, the angles marked with black dots are equal, as are the angles marked with red dots.  Therefore the total angles at the vertices are equal:  $\angle D = \angle C$.  

We have SAS, so $\triangle ABC \cong \triangle ABD$.

$\square$

To make the drawing, we have chosen to align along the longest side.  If we allow other sides to be used, then some configurations invalidate the proof.  But we can always choose the longest side, so the proof works for all triangles.  

SSS $\rightarrow$ SAS, and since SAS is sufficient to show congruence, so is SSS.

\subsection*{Prop. I.7}

\label{sec:Euclid7}

\begin{center} \includegraphics [scale=0.15] {Euclid_I_7c.png} \end{center}

$\bullet$  \ Let $\triangle ABC$ be drawn and a point $D$ be chosen on the same side of $AB$ as $C$ lies.  It cannot be true that both $AC = AD$ and $BC = BD$. 

\emph{Proof}.

Suppose that both statements are true: $AC = AD$ and $BC = BD$.

Then $\triangle ACD$ is isosceles so $\angle ACD = \angle ADC$.

We notice that 
\[ \angle BCD < \angle ACD = \angle ADC < \angle BDC \]

But since $\triangle BCD$ is also isosceles, $\angle BCD = \angle BDC$.

\begin{center} \includegraphics [scale=0.15] {Euclid_I_7c.png} \end{center}

This is a contradiction.  Therefore it cannot be that both $AC = AD$ and $BC = BD$.

$\square$

Going back to the perpendicular bisector

\begin{center} \includegraphics [scale=0.4] {iso13c.png} \end{center}

The forward theorem says that every point on the perpendicular bisector is equidistant from the two ends of $BC$.  The converse theorem says that every point which is equidistant, lies on the bisector.  No point not on the bisector can be equidistant.

\emph{Proof}

Let $P$ be not on the bisector.  Now, suppose that $PB = PC$.  Find $A$ as the point where either $PB$ or $PC$ cuts the bisector.  By the forward theorem, $AB = AC$.

But this is a contradiction, according to Euclid I.7.

$\square$

\subsection*{revisit Euclid I.7}
Perhaps you noticed, there is a problem with our proof of Euclid I.7.  Recall that $AC = AD$ and it is supposed that $BC = BD$.

\begin{center} \includegraphics [scale=0.13] {Euclid_I_7c.png} \end{center}
$D$ is drawn so that it is not contained within $\triangle ABC$.  But suppose it were?  

This illustrates a problem with drawing diagrams, we may introduce unrecognized assumptions into the proof.  The most common is to draw an acute triangle and presume it stands for all triangles.  

Here the problem is that the proof given before makes no sense if $D$ lies inside $\triangle ABC$.
\begin{center} \includegraphics [scale=0.15] {Euclid_I_7d.png} \end{center}

\emph{Proof} (Alternate).

We can find a different proof for this case in several ways.  Suppose $AC = AD$ and $BC = BD$.  Then $\triangle ABC \cong \triangle ABD$ by SSS, since $AB$ is shared.

But the base angles are \emph{not} equal (e.g. $\angle CAB > \angle DAB$).  This is a contradiction.

$\square$

This leads to a second difficulty, SSS is Euclid I.8, and Euclid actually proves I.8 by relying on I.7!

Luckily, we have a proof of SSS (given above) that does not depend on I.7, so this could still work.

Perhaps a better solution is this:  

\begin{center} \includegraphics [scale=0.15] {Euclid_I_7d.png} \end{center}

\emph{Proof}.

Suppose that both $AC = AD$ and $BC = BD$.

Then $\triangle ACD$ is isosceles so $\angle ACD = \angle ADC$.  Euclid I.5 also says that the supplementary angles are equal.  $\angle ECD = \angle FDC$.

We notice that 
\[ \angle BCD < \angle ECD = \angle FDC < \angle BDC \]

But since $\triangle BCD$ is also isosceles, $\angle BCD = \angle BDC$. 
This is a contradiction.  Therefore it cannot be that both $AC = AD$ and $BC = BD$.

$\square$

[ Another idea might be to use Euclid I.21, which says that if $D$ lies inside $\triangle ABC$, then $AD + BD < AC + BC$,  but that proposition turns out to be dependent through several steps, on I.8. ]

\subsection*{problems}

To prove:

$\circ$ \ Prove that for two supplementary angles, the angle bisectors are perpendicular to each other.

$\circ$ \ Prove that an equilateral triangle (all 3 sides equal) is equiangular (all 3 angles equal).  (Don't just rely on symmetry.  Adapt the proofs given in this chapter).

$\circ$ \ A line perpendicular to the bisector of an angle cuts off congruent segments on its sides.

In the figure below, given that $AC = AB$ and $\angle B = \angle C$.

\begin{center} \includegraphics [scale=0.4] {iso1.png} \end{center}

Prove that $BE = DC$.

Hint:  draw $BC$ and then mark all the angles that are equal.

\begin{center} \includegraphics [scale=0.4] {iso2.png} \end{center}

$\circ$ \ An equilateral triangle has all three angles equal.

\subsection*{problem}

\begin{center} \includegraphics [scale=0.4] {Hopkins_155.png} \end{center}

\begin{center} \includegraphics [scale=0.5] {iso_ext_prob.png} \end{center}

The angles labeled with black dots are equal by alternate interior angles and then vertical angles, while the angles labeled magenta are equal by alternate interior angles.

But we are given that this triangle is isosceles, so black and magenta are equal.  Therefore the exterior angle at $A$ is bisected by the horizontal.

Alternatively, use the fact that the exterior angle is the sum of the two base angles, and just use alternate interior angles once.

\end{document}